\documentclass[a4paper,12pt]{article}
\usepackage[utf8]{inputenc}
\usepackage[spanish]{babel}
\usepackage{tikz}
\usepackage{geometry}
\usepackage{setspace}
\usepackage{titlesec}
%\usepackage{lmodern}
\usepackage{multicol}
\usepackage{fancyhdr}
\usepackage{graphicx}
%\usepackage{helvet}
%\usepackage{mathptmx}
\usepackage{amsmath}       
\usepackage{amssymb}      
\usepackage{adjustbox}     
\usepackage{lastpage}
\usepackage{pgfplots}
\pgfplotsset{compat=1.17}
\usepackage{amsmath}
\usepackage{minted}
\usepackage{wrapfig}

\usepackage{tabularx}
\usepackage{qrcode}

\usepackage{hyperref}
\usepackage{booktabs}
\usepackage{array}
\usepackage{float}
\usepackage{longtable}  

\renewcommand{\arraystretch}{1.4}
\setlength{\tabcolsep}{8pt} 

\setlength{\footskip}{25pt}
%---------------------------python
\usepackage[T1]{fontenc}
\usepackage[utf8]{inputenc}

\usepackage{listings}
\usepackage{xcolor}

\lstdefinestyle{pythonstyle}{
	language=Python,
	basicstyle=\ttfamily\small,
	keywordstyle=\color{blue},
	stringstyle=\color{red},
	commentstyle=\color{gray},
	numbers=left,
	numberstyle=\tiny\color{gray},
	stepnumber=1,
	numbersep=5pt,
	showspaces=false,
	showstringspaces=false,
	breaklines=true,
	frame=single,
	tabsize=4
}



%-------------------------------Consola
\usepackage{listings}
\usepackage{xcolor}

\lstdefinestyle{consola}{
	backgroundcolor=\color{black!5},
	basicstyle=\ttfamily,
	frame=single,
	keywordstyle=\color{blue},
	commentstyle=\color{gray},
	showstringspaces=false
}

%Plano Z y Z^2

\usepackage[utf8]{inputenc}
\usepackage{amsmath, amssymb}
\usepackage{pgfplots}
\pgfplotsset{compat=1.18}
\usepackage{geometry}
\geometry{a4paper, margin=2.5cm}


\geometry{top=2.5cm, bottom=2.5cm, left=2.5cm, right=1.5cm}

% Estilo de página
\pagestyle{fancy}
\fancypagestyle{miestilo}{
	\fancyhf{} % Limpia encabezado y pie
	
	% Encabezado izquierda: logo
	\fancyhead[L]{%
		\vspace{0.1cm}
		\adjustbox{valign=c,scale=0.15}{%
			\begin{tikzpicture}
				\clip (-3,-2.8) rectangle (3,2.8);
				\draw[black,line width=0.8cm] (-2.3,0)--(2.3,0);
				\draw[black,line width=0.8cm] (0,-2.85)--(0,2.85);
				\draw[black,line width=0.87cm] (0,2.8) circle(2);
				\draw[black,line width=0.87cm] (0,-2.8) circle(2);
			\end{tikzpicture}
		}%
	}
	
	% Encabezado centro: título
	\fancyhead[C]{\textbf{Dispositivos Electrónicos}}
	
	% Encabezado derecha: TP y año debajo
	\fancyhead[R]{%
		Trabajo Practico de Laboratorio 1 \\ 2026
	}
	
	% Línea en el encabezado y pie
	\renewcommand{\headrulewidth}{0.4pt}
	\renewcommand{\footrulewidth}{0.4pt}
	
	% Pie de página derecha: Página X de Y
	\fancyfoot[R]{Página \thepage\ de \pageref{LastPage}}
}


\begin{document}
	
	% --------- PORTADA ---------
	\pagestyle{empty}
	\begin{center}
		\textbf{\LARGE UNIVERSIDAD TECNOLÓGICA NACIONAL}\\[0.2cm]
		\textbf{\Large FACULTAD REGIONAL CÓRDOBA}\\[0.5cm]
		\textbf{\large INGENIERÍA ELECTRÓNICA}\\[1.5cm]
		
		\begin{tikzpicture}
			\clip (-3,-2.8) rectangle (3,2.8);
			\draw[black,line width=0.8cm] (-2.3,0)--(2.3,0);
			\draw[black,line width=0.8cm] (0,-2.85)--(0,2.85);
			\draw[black,line width=0.87cm] (0,2.8) circle(2);
			\draw[black,line width=0.87cm] (0,-2.8) circle(2);
		\end{tikzpicture}
		
		\vspace{1cm}
		\textbf{\huge Dispositivos Electrónicos }\\[0.5cm]
		\rule{\textwidth}{0.4pt}\\[0.3cm]
		\textbf{\Huge Trabajo Practico de Laboratorio 1}\\[0.3cm]
		{\large Identificación de Componentes Electrónicos y
			Mediciones Básicas}\\[0.3cm]
		\rule{\textwidth}{0.4pt}\\[1.5cm]
	\end{center}
	
	\begin{flushleft}
		\begin{tabular}{@{}ll}
			DOCENTES     & : Ing. Luis Alberto Guanuco\\ 
			& \ \ Ing. Leonardo Gabriel Gamboa \\[0.2cm]
			CÁTEDRA      & :  Análisis de Señales y Sistemas  \\[0.2cm]
			CURSO        & : 3R1 \\[0.2cm]
			ALUMNO      & :
			\begin{tabular}[t]{@{}l l}
				Gonza Elías Noel Imanol           & 87669 \\ 
				
				Juarez Luciano                    & 95021   \\ 
				Pereyra Nahuel                    & 402333  \\ 
				         
				
			\end{tabular}
		\end{tabular}
	\end{flushleft}
	
	\vfill
	\begin{center}
		\textbf{CÓRDOBA, ARGENTINA}\\
		\textbf{2026}
	\end{center}
	\pagestyle{empty}
	% --------- CAMBIO DE ESTILO ---------
	\newpage
	\pagenumbering{arabic}
	\pagestyle{miestilo}
	% --------- CONTENIDO ---------
	\tableofcontents
	\newpage
	
	
	\section*{Introducción}
	\addcontentsline{toc}{section}{Introducción}
	En el presente Trabajo Práctico se abordó el estudio de circuitos resistivos de corriente continua, con el objetivo de aplicar los conceptos teóricos fundamentales de análisis eléctrico y contrastarlos con resultados obtenidos tanto en simulación como en laboratorio.
	
	En una primera etapa, se realizó el análisis del circuito propuesto utilizando la Ley de Ohm y las Leyes de Kirchhoff, determinando las tensiones, corrientes y potencias presentes en cada uno de los elementos. Esta instancia permitió comprender el comportamiento eléctrico del circuito a partir de expresiones teóricas y resoluciones analíticas.
	
	Posteriormente, se empleó el software LTspice para simular el circuito y obtener sus valores de operación, verificando la coherencia entre los resultados calculados y los arrojados por la simulación. Además, se analizaron distintas herramientas del simulador para observar el comportamiento del circuito frente a variaciones de sus parámetros.
	
	Finalmente, se implementó físicamente el circuito en protoboard y se realizaron mediciones experimentales mediante el uso de multímetro, registrando los valores de tensión, corriente y potencia. A partir de la comparación entre análisis teórico, simulación e implementación práctica, se evaluaron las diferencias observadas y se identificaron las posibles causas asociadas a tolerancias de los componentes, condiciones reales de medición y limitaciones de los modelos ideales.
	
	\newpage
	
	
	
	\section*{Análisis del Circuito}
	\addcontentsline{toc}{section}{Análisis del Circuito}
	
	\begin{itemize}
		\item $V_S = 10[V]$
		\item $R_1 = 10[k\Omega]$
		\item $R_2 = 4,7[k\Omega]$
		\item $R_3 = 3,3[k\Omega]$
	\end{itemize}
	
	\begin{figure}[H]
		\centering
		\includegraphics[width=0.45\textwidth]{imagenes/img2.png}
		\caption{Circuito para análisis.}
		\label{fig:circuito_analisis}
	\end{figure}
	
	
	Para comenzar con el análisis del circuito, se realizó una revisión de los conceptos principales de la
	Ley de Ohm y de las Leyes de Kirchhoff. A partir de esto, se identificó que el circuito presentado
	corresponde a un circuito mixto, ya que posee una resistencia conectada en serie y dos resistencias
	conectadas en paralelo.
	
	Este reconocimiento es importante, ya que permite aplicar correctamente los métodos de resolución
	para calcular la resistencia total del circuito, la corriente total y las corrientes que circulan por cada
	resistencia.
	
	
	
	
	
	
	
	
\begin{figure}[H]
	\centering
	
	\begin{minipage}{0.30\textwidth}
		\centering
		\includegraphics[width=\textwidth]{imagenes/img3.png}
		\caption{Circuito equivalente reemplazando el paralelo por $R_{23}$.}
		\label{fig:circuito_equivalente}
	\end{minipage}
	\hfill
	\begin{minipage}{0.60\textwidth}
		En una conexión serie, la corriente es la misma en todos los elementos. Por esta razón, la corriente
		que circula por $R_1$ es igual a la corriente total. Sin embargo, las caídas de tensión en cada resistencia
		pueden ser diferentes, ya que dependen del valor de cada resistencia.
		
		Luego de obtener la corriente total, se continúa con el cálculo de las corrientes $I_2$ e $I_3$, que
		corresponden a las corrientes que circulan por las resistencias conectadas en paralelo.
		
		Por lo tanto:
		
		\[
		R_{23} = 1938,75\Omega
		\]
		
		\[
		I_1 = I_T = 0,8376\,mA
		\]
	\end{minipage}
	
\end{figure}
	
	
	
	
	\subsection*{Aplicación de las Leyes de Kirchhoff}
	\addcontentsline{toc}{subsection}{Aplicación de las Leyes de Kirchhoff}
	
	Para continuar con el análisis del circuito, se aplican las Leyes de Kirchhoff.
	
	\subsubsection*{Ley de Corrientes de Kirchhoff}
	\addcontentsline{toc}{subsubsection}{Ley de Corrientes de Kirchhoff}
	
	Según la Ley de Corrientes de Kirchhoff, la corriente total que llega al nodo se divide entre las
	ramas en paralelo. Por lo tanto:
	
	\[
	I_T = I_2 + I_3
	\]
	
	\subsubsection*{Ley de Tensiones de Kirchhoff}
	\addcontentsline{toc}{subsubsection}{Ley de Tensiones de Kirchhoff}
	
	Aplicando la Ley de Tensiones de Kirchhoff en la malla principal, se obtiene:
	
	\[
	V_S - V_{R1} - V_{R2} = 0
	\]
	
	Por lo tanto:
	
	\[
	V_S = V_{R1} + V_{R2}
	\]
	
	Como las caídas de tensión se calculan mediante la Ley de Ohm:
	
	\[
	V = I \cdot R
	\]
	
	entonces:
	
	\[
	V_S = I_T \cdot R_1 + I_2 \cdot R_2
	\]
	
	Para la segunda malla, correspondiente al paralelo entre $R_2$ y $R_3$, se cumple que la caída de
	tensión en ambas resistencias es la misma:
	
	\[
	V_{R2} - V_{R3} = 0
	\]
	
	\[
	V_{R2} = V_{R3}
	\]
	
	Reemplazando por Ley de Ohm:
	
	\[
	I_2 \cdot R_2 = I_3 \cdot R_3
	\]
	
	Despejando $I_3$:
	
	\[
	I_3 = I_2 \cdot \frac{R_2}{R_3}
	\]
	
	Luego, reemplazamos esta expresión en la Ley de Corrientes de Kirchhoff:
	
	\[
	I_T = I_2 + I_3
	\]
	
	\[
	I_T = I_2 + I_2 \cdot \frac{R_2}{R_3}
	\]
	
	Sacando factor común $I_2$:
	
	\[
	I_T = I_2 \left(1 + \frac{R_2}{R_3}\right)
	\]
	
	\[
	I_T = I_2 \left(\frac{R_3}{R_3} + \frac{R_2}{R_3}\right)
	\]
	
	\[
	I_T = I_2 \left(\frac{R_3 + R_2}{R_3}\right)
	\]
	
	Despejando $I_2$:
	
	\[
	I_2 = I_T \left(\frac{R_3}{R_2 + R_3}\right)
	\]
	
	Y si querés agregar después los valores numéricos de las caídas de tensión:
	
	\subsubsection*{Cálculo de caídas de tensión}
	\addcontentsline{toc}{subsubsection}{Cálculo de caídas de tensión}
	
	La caída de tensión en $R_1$ se calcula como:
	
	\[
	V_{R1} = I_T \cdot R_1
	\]
	
	\[
	V_{R1} = 0,837\,mA \cdot 10k\Omega
	\]
	
	\[
	V_{R1} = 8,37V
	\]
	
	La caída de tensión en $R_2$ es:
	
	\[
	V_{R2} = I_2 \cdot R_2
	\]
	
	\[
	V_{R2} = 0,346\,mA \cdot 4,7k\Omega
	\]
	
	\[
	V_{R2} = 1,62V
	\]
	
	La caída de tensión en $R_3$ es:
	
	\[
	V_{R3} = I_3 \cdot R_3
	\]
	
	\[
	V_{R3} = 0,491\,mA \cdot 3,3k\Omega
	\]
	
	\[
	V_{R3} = 1,62V
	\]
	
	Como $R_2$ y $R_3$ están conectadas en paralelo, ambas tienen la misma caída de tensión:
	
	\[
	V_{R2} = V_{R3} = 1,62V
	\]
	
	
	
	
	\subsubsection*{Despeje de la corriente total}
	\addcontentsline{toc}{subsubsection}{Despeje de la corriente total}
	
	
	\begin{figure}[H]
		\centering
		
		\begin{minipage}{0.30\textwidth}
			\centering
			\includegraphics[width=\textwidth]{imagenes/img3.png}
			\caption{Circuito equivalente reemplazando el paralelo por $R_{23}$.}
			\label{fig:circuito_equivalente}
		\end{minipage}
		\hfill
		\begin{minipage}{0.60\textwidth}
		
		Partiendo de la ecuación de la malla principal:
		
		\[
		V_S = I_T R_1 + I_2 R_2
		\]
		
		Como anteriormente se obtuvo que:
		
		\[
		I_2 = I_T \left(\frac{R_3}{R_2 + R_3}\right)
		\]
		
		se reemplaza $I_2$ en la ecuación de tensión:
		
		\[
		V_S = I_T R_1 + I_T \left(\frac{R_3}{R_2 + R_3}\right) R_2
		\]
		
		Luego, se extrae factor común $I_T$:
		
		\[
		V_S = I_T \left( R_1 + \frac{R_3 R_2}{R_2 + R_3} \right)
		\]
		
		Finalmente, se despeja la corriente total:
		
		\[
		I_T =
		\frac{V_S}
		{R_1 + \frac{R_3 R_2}{R_2 + R_3}}
		\]
		
		Esta expresión permite calcular la corriente total del circuito considerando la resistencia $R_1$ en serie
		con el paralelo formado por $R_2$ y $R_3$.
		

	\end{minipage}
	
	\end{figure}
	
	
	
	
	

	
	
	\subsubsection*{Cálculo de la resistencia equivalente}
	\addcontentsline{toc}{subsubsection}{Cálculo de la resistencia equivalente}
	
	Como $R_2$ y $R_3$ se encuentran conectadas en paralelo, primero se calcula la resistencia equivalente
	de esa parte del circuito:
	
	\[
	R_{eq} = R_1 + (R_2 \parallel R_3)
	\]
	
	\[
	R_{23} = \frac{R_2 \cdot R_3}{R_2 + R_3}
	\]
	
	Reemplazando los valores:
	
	\[
	R_{23} =
	\frac{4,7k\Omega \cdot 3,3k\Omega}{4,7k\Omega + 3,3k\Omega}
	\]
	
	\[
	R_{23} =
	\frac{15,51k\Omega^2}{8k\Omega}
	\]
	
	\[
	R_{23} = 1,94k\Omega
	\]
	
	Luego, como $R_1$ está en serie con $R_{23}$, se calcula la resistencia total:
	
	\[
	R_T = R_1 + R_{23}
	\]
	
	\[
	R_T = 10k\Omega + 1,94k\Omega
	\]
	
	\[
	R_T = 11,94k\Omega
	\]
	
	
	\subsubsection*{Cálculo de las corrientes del circuito}
	\addcontentsline{toc}{subsubsection}{Cálculo de las corrientes del circuito}
	
	A partir de la resistencia total del circuito, se calcula la corriente total mediante la Ley de Ohm:
	
	\[
	I_T =
	\frac{V_S}
	{R_1 + \frac{R_3 R_2}{R_3 + R_2}}
	\]
	
	Reemplazando los valores:
	
	\[
	I_T =
	\frac{10V}
	{10k\Omega + \frac{3,3k\Omega \cdot 4,7k\Omega}{3,3k\Omega + 4,7k\Omega}}
	\]
	
	\[
	I_T =
	\frac{10V}
	{10k\Omega + \frac{15,51k\Omega^2}{8k\Omega}}
	\]
	
	\[
	I_T =
	\frac{10V}{10k\Omega + 1,938k\Omega}
	\]
	
	\[
	I_T =
	\frac{10V}{11,938k\Omega}
	\]
	
	\[
	I_T = 0,837mA
	\]
	
	Como $R_1$ se encuentra en serie con el resto del circuito, la corriente que circula por $R_1$ es igual
	a la corriente total:
	
	\[
	I_1 = I_T
	\]
	
	\[
	I_1 = 0,837mA
	\]
	
	Luego, para calcular la corriente que circula por $R_2$, se utiliza la relación de división de corriente:
	
	\[
	I_2 = I_T \left(\frac{R_3}{R_3 + R_2}\right)
	\]
	
	Reemplazando los valores:
	
	\[
	I_2 =
	0,837mA
	\left(
	\frac{3,3k\Omega}{3,3k\Omega + 4,7k\Omega}
	\right)
	\]
	
	\[
	I_2 =
	0,837mA
	\left(
	\frac{3,3k\Omega}{8k\Omega}
	\right)
	\]
	
	\[
	I_2 =
	0,837mA \cdot 0,413
	\]
	
	\[
	I_2 = 0,346mA
	\]
	
	Finalmente, aplicando la Ley de Corrientes de Kirchhoff:
	
	\[
	I_T = I_2 + I_3
	\]
	
	se despeja $I_3$:
	
	\[
	I_3 = I_T - I_2
	\]
	
	\[
	I_3 = 0,837mA - 0,346mA
	\]
	
	\[
	I_3 = 0,491mA
	\]
	
	
	
	
	
	\subsubsection*{Cálculo de potencia}
	\addcontentsline{toc}{subsubsection}{Cálculo de potencia}
	
	Una vez obtenidos los valores de tensión y corriente en cada resistencia, se calcula la potencia
	utilizando la expresión:
	
	\[
	P = V \cdot I
	\]
	
	Los valores obtenidos se resumen en la siguiente tabla:
	
	\begin{table}[H]
		\centering
		\begin{tabular}{|c|c|c|c|}
			\hline
			\textbf{Elemento} & \textbf{Tensión (V)} & \textbf{Corriente (mA)} & \textbf{Potencia (mW)} \\
			\hline
			$V_S$ & $10V$ & $0,837mA$ & $8,37mW$ \\
			\hline
			$R_1$ & $8,37V$ & $0,837mA$ & $7,01mW$ \\
			\hline
			$R_2$ & $1,62V$ & $0,346mA$ & $0,56mW$ \\
			\hline
			$R_3$ & $1,62V$ & $0,491mA$ & $0,80mW$ \\
			\hline
		\end{tabular}
		\caption{Valores de tensión, corriente y potencia del circuito.}
		\label{tab:valores_potencia}
	\end{table}
	
	A continuación, se muestran los cálculos de potencia para cada elemento del circuito:
	
	\[
	P_{R1} = V_{R1} \cdot I_T
	\]
	
	\[
	P_{R1} = 8,37V \cdot 0,837mA
	\]
	
	\[
	P_{R1} = 7,01mW
	\]
	
	\[
	P_{R2} = V_{R2} \cdot I_2
	\]
	
	\[
	P_{R2} = 1,62V \cdot 0,346mA
	\]
	
	\[
	P_{R2} = 0,56mW
	\]
	
	\[
	P_{R3} = V_{R3} \cdot I_3
	\]
	
	\[
	P_{R3} = 1,62V \cdot 0,491mA
	\]
	
	\[
	P_{R3} = 0,80mW
	\]
	
	Finalmente, la potencia total entregada por la fuente se calcula como:
	
	\[
	P_S = V_S \cdot I_T
	\]
	
	\[
	P_S = 10V \cdot 0,837mA
	\]
	
	\[
	P_S = 8,37mW
	\]
	
	
	
   \section*{Simulación en LTspice}
		\addcontentsline{toc}{section}{Simulación en LTspice}
	

	\subsection*{Simulación DC op point}
	\addcontentsline{toc}{subsection}{Simulación DC op point}
	
	
	
	Se armó el circuito en LTspice y se ejecutó el análisis de punto de operación DC. En la fuente y en
	algunas corrientes puede aparecer signo negativo por la convención de sentido que toma LTspice; 
	físicamente indica que la fuente entrega energía al circuito.
	
	\begin{figure}[H]
		\centering
		\includegraphics[width=0.85\textwidth]{imagenes/img1.png}
		\caption{Circuito simulado y resultados de punto de operación en LTspice.}
		\label{fig:simulacion_ltspice_op}
	\end{figure}
	
	\begin{table}[H]
		\centering
		\renewcommand{\arraystretch}{1.2}
		\begin{tabular}{|c|c|c|c|}
			\hline
			
			\textbf{Elemento} & \textbf{Tensión (V)} & \textbf{Corriente (mA)} & \textbf{Potencia (mW)} \\
			\hline
			$V_S$ & $10,000$ & $0,838$ & $8,376$ \\
			\hline
			$R_1$ & $8,376$ & $0,838$ & $7,016$ \\
			\hline
			$R_2$ & $1,624$ & $0,346$ & $0,561$ \\
			\hline
			$R_3$ & $1,624$ & $0,492$ & $0,799$ \\
			\hline
		\end{tabular}
		\caption{Valores esperados por simulación DC op point.}
		\label{tab:valores_simulacion_dc}
	\end{table}
	
	\subsection*{Barrido de tensión continua}
	\addcontentsline{toc}{subsection}{Barrido de tensión continua}
	
	
	En el barrido DC se varió la fuente $V_S$ desde $0\,V$ hasta $20\,V$. Las corrientes resultaron
	proporcionales a la tensión aplicada, tal como predice la Ley de Ohm para una red resistiva lineal.
	
	\begin{figure}[H]
		\centering
		\includegraphics[width=0.6\textwidth]{imagenes/img4.png}
		\caption{Barrido DC de corriente en $R_1$, $R_2$ y $R_3$ al variar $V_S$.}
		\label{fig:barrido_dc_corrientes}
	\end{figure}
	
	
	
	\subsection*{Barrido de parámetro $R_2$}
	\addcontentsline{toc}{subsection}{Barrido de parámetro $R_2$}
	
	Se realizó un barrido paramétrico de $R_2$ desde $1\,k\Omega$ hasta $10\,k\Omega$.
	Al aumentar $R_2$, aumenta la tensión del nodo B porque la resistencia equivalente del paralelo
	crece. Al mismo tiempo, la corriente por $R_2$ disminuye porque su resistencia aumenta.
	
	\begin{figure}[H]
		\centering
		\includegraphics[width=0.60\textwidth]{imagenes/img5.png}
		\caption{Barrido paramétrico de $R_2$ y variación de $V(B)$ e $I(R_2)$.}
		\label{fig:barrido_parametrico_r2}
	\end{figure}
	
	\begin{table}[H]
		\centering
		\renewcommand{\arraystretch}{1.2}
		\begin{tabular}{|c|c|c|c|}
			\hline
		
			\textbf{$R_2$ ($k\Omega$)} & \textbf{$V(B)$ (V)} & \textbf{$I(R_2)$ (mA)} & \textbf{$I(R_3)$ (mA)} \\
			\hline
			$1,0$  & $0,713$ & $0,713$ & $0,216$ \\
			\hline
			$2,0$  & $1,107$ & $0,554$ & $0,336$ \\
			\hline
			$3,0$  & $1,358$ & $0,453$ & $0,412$ \\
			\hline
			$4,7$  & $1,624$ & $0,346$ & $0,492$ \\
			\hline
			$10,0$ & $1,988$ & $0,199$ & $0,602$ \\
			\hline
		\end{tabular}
		\caption{Valores representativos del barrido paramétrico de $R_2$.}
		\label{tab:barrido_parametrico_r2}
	\end{table}
	
	
	\section*{Implementación y mediciones}
		\addcontentsline{toc}{section}{Implementación y mediciones}
	
	Para la implementación se propone armar el circuito en protoboard con resistores axiales de
	$\frac{1}{4}\,W$  y alimentar el circuito con $V_S = 10\,V$.
	Las mediciones de tensión se realizan en paralelo con cada componente y las de corriente en serie,
	o de manera indirecta mediante la relación:
	
	\[
	I = \frac{V}{R}
	\]
	
	\begin{table}[H]
		\centering
		\renewcommand{\arraystretch}{1.2}
		\begin{tabular}{|c|c|c|}
			\hline
		
			\textbf{Resistor} & \textbf{Valor nominal} & \textbf{Valor usado en el cálculo} \\
			\hline
			$R_1$ & $10\,k\Omega$  & $9,86\,k\Omega$ \\
			\hline
			$R_2$ & $4,7\,k\Omega$ & $4,622\,k\Omega$ \\
			\hline
			$R_3$ & $3,3\,k\Omega$ & $3,21\,k\Omega$ \\
			\hline
		\end{tabular}
		\caption{Valores de resistencias utilizadas.}
		\label{tab:resistencias_utilizadas}
	\end{table}
	\begin{table}[H]
		\centering
		\renewcommand{\arraystretch}{1.2}
		
		\resizebox{\textwidth}{!}{
			\begin{tabular}{|c|c|c|c|}
				\hline
				\textbf{Elemento} &
				\textbf{Tensión medida/esperada (V)} &
				\textbf{Corriente medida/esperada (mA)} &
				\textbf{Potencia medida/esperada (mW)} \\
				\hline
				$V_S$ & $10,000$ & $0,837$ & $8,370$ \\
				\hline
				$R_1$ & $8,370$ & $0,837$ & $7,010$ \\
				\hline
				$R_2$ & $1,620$ & $0,346$ & $0,560$ \\
				\hline
				$R_3$ & $1,620$ & $0,491$ & $0,800$ \\
				\hline
			\end{tabular}
		}
		
		\caption{Valores de implementación organizados según los datos aportados.}
		\label{tab:valores_implementacion}
	\end{table}
	
	
	
	\subsection*{Comparación de resultados}
			\addcontentsline{toc}{subsection}{Comparación de resultados}
	
	\begin{table}[H]
		\centering
		\small
		\renewcommand{\arraystretch}{1.2}
		\resizebox{\textwidth}{!}{
			\begin{tabular}{|c|c|c|c|c|c|}
				\hline
				
				\textbf{Magnitud} & \textbf{Origen} & \textbf{$V_S$} & \textbf{$R_1$} & \textbf{$R_2$} & \textbf{$R_3$} \\
				\hline
				Tensión (V) & Análisis & $10,000$ & $8,376$ & $1,624$ & $1,624$ \\
				\hline
				Tensión (V) & Simulación & $10,000$ & $8,376$ & $1,624$ & $1,624$ \\
				\hline
				Tensión (V) & Medición/esperado & $10,000$ & $8,370$ & $1,620$ & $1,620$ \\
				\hline
				Corriente (mA) & Análisis & $0,838$ & $0,838$ & $0,346$ & $0,492$ \\
				\hline
				Corriente (mA) & Simulación & $0,838$ & $0,838$ & $0,346$ & $0,492$ \\
				\hline
				Corriente (mA) & Medición/esperado & $0,837$ & $0,837$ & $0,346$ & $0,491$ \\
				\hline
				Potencia (mW) & Análisis & $8,376$ & $7,016$ & $0,561$ & $0,799$ \\
				\hline
				Potencia (mW) & Simulación & $8,376$ & $7,016$ & $0,561$ & $0,799$ \\
				\hline
				Potencia (mW) & Medición/esperado & $8,370$ & $7,010$ & $0,560$ & $0,800$ \\
				\hline
			\end{tabular}
		}
		\caption{Comparación entre análisis, simulación y medición.}
		\label{tab:comparacion_resultados}
	\end{table}
	
	Los valores obtenidos mediante el análisis teórico y la simulación coinciden dentro del error de
	redondeo. Las pequeñas diferencias esperadas en laboratorio se deben principalmente a la tolerancia
	real de los resistores, a la precisión del multímetro, a la resistencia interna del instrumento, a los
	contactos de la protoboard y al ajuste exacto de la fuente de alimentación.
	
	
	\section*{Conclusiones}		
	\addcontentsline{toc}{section}{Conclusiones}
	
	
	A partir del análisis realizado, se puede concluir que el circuito estudiado corresponde a una red
	mixta, formada por una resistencia en serie con un conjunto de resistencias en paralelo. Para resolverlo,
	primero se calculó la resistencia equivalente del paralelo entre $R_2$ y $R_3$, y luego se sumó este
	resultado con $R_1$ para obtener la resistencia total del circuito.
	
	Los resultados obtenidos mediante el cálculo teórico fueron coherentes con los valores simulados en
	LTspice. Esto permite verificar que la aplicación de la Ley de Ohm y de las Leyes de Kirchhoff fue
	correcta. Además, se comprobó que la corriente total que entrega la fuente es la misma que circula por
	$R_1$, mientras que en el nodo B dicha corriente se divide entre las ramas de $R_2$ y $R_3$.
	
	También se verificó que las resistencias conectadas en paralelo poseen la misma caída de tensión. En
	este caso, tanto $R_2$ como $R_3$ presentan aproximadamente:
	
	\[
	V_{R2} = V_{R3} = 1,62\,V
	\]
	
	Por otro lado, se calculó la potencia disipada en cada elemento. Se observó que las potencias máximas
	son muy inferiores a $1/4\,W$ y a $1/8\,W$. La resistencia más exigida es $R_1$, con una potencia
	aproximada de:
	
	\[
	P_{R1} = 7,016\,mW
	\]
	
	Por lo tanto, los resistores seleccionados trabajan con un amplio margen de seguridad. Comparando
	con resistores de $1/4\,W$, se tiene aproximadamente:
	
	\[
	R_1 = 2,81\%, \qquad R_2 = 0,22\%, \qquad R_3 = 0,32\%
	\]
	
	Y para resistores de $1/8\,W$:
	
	\[
	R_1 = 5,61\%, \qquad R_2 = 0,45\%, \qquad R_3 = 0,64\%
	\]
	
	Esto demuestra que no existe riesgo de sobrecalentamiento en las resistencias bajo las condiciones
	analizadas.
	
	En la simulación ideal de LTspice no siempre se consideran factores reales como la tolerancia de los
	resistores, la variación por temperatura, la resistencia de cables y contactos, la resistencia interna de
	la fuente, la resistencia interna del multímetro y posibles falsos contactos en la protoboard. Por ese
	motivo, al realizar mediciones prácticas pueden aparecer pequeñas diferencias respecto de los valores
	teóricos y simulados.
	
	El barrido paramétrico de $R_2$ permitió observar que, al modificar esta resistencia, también cambia
	la resistencia equivalente del paralelo. Cuando $R_2$ aumenta, la tensión del nodo B aumenta y la
	corriente que circula por $R_2$ disminuye. Por lo tanto, $R_2$ afecta el reparto de corriente entre las
	ramas paralelas.
	
	Finalmente, el análisis permitió comprobar que el cálculo teórico, la simulación y la medición son
	herramientas complementarias. El cálculo permite predecir el comportamiento del circuito, la simulación
	ayuda a verificarlo antes de armarlo físicamente y la medición permite confirmar el funcionamiento real.
	
	
	
	\subsection*{Selección de potencia nominal}
	\addcontentsline{toc}{subsection}{Selección de potencia nominal}
	
	Las potencias máximas calculadas son muy inferiores a $1/4\,W$ y a $1/8\,W$.
	El resistor más exigido es $R_1$, con aproximadamente $7,016\,mW$, por lo que
	los resistores seleccionados trabajan con amplio margen de seguridad.
	
	\subsection*{Porcentaje de capacidad}
	\addcontentsline{toc}{subsection}{Porcentaje de capacidad}
	
	\[
	\text{Porcentaje de trabajo} =
	\frac{P_{\text{disipada}}}{P_{\text{nominal}}} \cdot 100
	\]
	
	Si se utilizan resistores de $1/4\,W$, los porcentajes de trabajo son:
	
	\[
	R_1 = 2,81\%, \qquad R_2 = 0,22\%, \qquad R_3 = 0,32\%
	\]
	
	Si se utilizan resistores de $1/8\,W$, los porcentajes de trabajo son:
	
	\[
	R_1 = 5,61\%, \qquad R_2 = 0,45\%, \qquad R_3 = 0,64\%
	\]
	
	Estos porcentajes son bajos y favorecen una buena vida útil de los componentes.
	
	\subsection*{Factores reales no contemplados por LTspice}
	\addcontentsline{toc}{subsection}{Factores reales no contemplados por LTspice}
	
	En la simulación ideal no siempre se consideran la tolerancia de los resistores, la variación por
	temperatura, la resistencia de cables y contactos, la resistencia interna de la fuente, la resistencia
	interna del multímetro ni posibles falsos contactos en la protoboard.
	
	\subsection*{Efecto del barrido de $R_2$}
	\addcontentsline{toc}{subsection}{Efecto del barrido de R2}
	
	La variación de $R_2$ modifica la resistencia equivalente del paralelo. Cuando $R_2$ aumenta, sube
	la tensión del nodo B y disminuye la corriente por $R_2$; por lo tanto, $R_2$ afecta el reparto de
	corriente en las ramas paralelas.
	
	\subsection*{Influencia de tolerancia}
	\addcontentsline{toc}{subsection}{Influencia de tolerancia}
	
	Si el valor real de una resistencia se aparta de su valor nominal, también se modifican las corrientes
	y tensiones medidas. Por eso puede aparecer una desviación entre los valores teóricos y los medidos.
	
	\subsection*{Efecto del amperímetro}
	\addcontentsline{toc}{subsection}{Efecto del amperímetro}
	
	Al conectar el multímetro en serie para medir corriente se incorpora una pequeña resistencia interna
	que puede modificar el circuito. En este caso, al trabajar con resistencias de kiloohmios, el efecto
	suele ser pequeño, pero no debe ignorarse.
	
	\subsection*{Método más fiable}
	\addcontentsline{toc}{subsection}{Método más fiable}
	
	Para la toma de decisiones de ingeniería, la medición experimental es la validación final del
	comportamiento real. Sin embargo, el cálculo y la simulación son indispensables para prever
	magnitudes, detectar errores y seleccionar componentes antes de construir.
	
	\subsection*{Variaciones entre análisis, simulación y medición}
	\addcontentsline{toc}{subsection}{Variaciones entre análisis, simulación y medición}
	
	Las variaciones observadas son reducidas. Las principales causas son redondeo, tolerancia de
	componentes, precisión instrumental, ajuste de la fuente y condiciones reales de montaje.
	
	
\end{document}